%! Describing the Distribution of a Quantitative Variable — Unit 1, Lesson 1.4 — Experience & Formalize
%! (Activity · QuickNotes · Application)
% EFFL component, Math Medic sizing: 12pt body, OPEN white space after every question.
% Page budget: Activity <=2pp · QuickNotes 1/2pp · Application 1/2-1pp. TOTAL 3pp.
% Check Your Understanding is NOT in this component: those items were moved to the end of
% the packet and are now the lesson's SCORED homework (see ../homework/).
%
% PAGE LOCKSTEP + PAGE BUDGET: the Activity reserves 10.8cm of \answerspace across its six
% written sub-questions (1.6 + 1.6 + 2.0 + 1.6 + 1.8 + 2.2), which is the same total the model
% lesson (unit01/lesson03) reserves and what holds this component to three pages. Do not raise
% a height here without lowering another, and never raise one in only one of the two files.
%
% SPOILER RULE: the Activity never says shape, symmetric, skewed, unimodal, bimodal, uniform,
% center, variability, spread, unusual feature, outlier, gap, cluster, or SOCS. It says
% "the middle wait", "bunched", "stragglers", "sits apart", "an empty stretch".
% `Distribution', `dotplot', and `histogram' WERE formalized in Lessons 1.1 and 1.3 and are
% prior knowledge here.
%
% NO SKETCH QUESTIONS (hard constraint): every display is pre-drawn. Students read, compare,
% and describe — they never draw an axis or place a dot.
%
% THE DATA
%   Line A (20 waits, min): 3x4  4x8  5x5  6x3        -> median 4, mean 4.35, range 3-6
%   Line B (20 waits, min): 1x4  2x3  3x2  4x2  5x3  6x2  7x2  11  13
%                                                     -> median 4, mean 4.45, range 1-13
%   Identical middle wait, completely different experience. That is the discovery.
%   Picture I  home prices, 40 homes, $100k bins:  6 14 10 5 3 1 1   (long right tail)
%   Picture II quiz scores, 32 students, 10-pt bins: 1 2 5 11 13     (long left tail)
%   Picture III club heights, 40 adults, 2-in bins: 2 5 9 4 0 0 4 9 5 2
%                       -> symmetric, two humps, empty stretch 66-70 in. A single "typical
%                          height" of about 67 in belongs to nobody in the club. That is the crux.
\documentclass[12pt]{article}
\usepackage{apstats-article}
\usepackage{apstats-boxes}
\newcommand{\answerspace}[2]{\par\nopagebreak\noindent\begin{minipage}[t][#1][t]{\linewidth}%
  \color{keyred}\bfseries #2\end{minipage}\par}

\begin{document}

\pageheader{Unit 1, Lesson 1.4}{Experience \& Formalize: Two Checkout Lines}
% NAMESTRIP: no \namedateperiod here — the name row lives on the cover only.

\vspace{0.02in}

% ── Part 1: the Activity ─────────────────────────────────────────────────────
\begin{headlinebox}{sky}
\textbf{Two Checkout Lines.} A supermarket timed 20 customers in each of its two checkout
lines. Both lines have \emph{exactly the same} middle wait. Work the two problems with your
group using only what you already know.
\end{headlinebox}

\vspace{0.02in}

\begin{scenariobox}[1. Which Line Would You Join?]{navy}
Each dot is one customer. Both plots use the same scale, so you can compare them directly.

\vspace{4pt}
\begin{center}
\begin{tikzpicture}[scale=0.80]
  % ── Line A: tightly bunched ──
  \draw[->, thick] (-0.2,0) -- (6.0,0);
  \foreach \x/\lab in {0/0, 0.7/2, 1.4/4, 2.1/6, 2.8/8, 3.5/10, 4.2/12, 4.9/14, 5.6/16} {
    \draw (\x,-0.07) -- (\x,0.07); \node[below, font=\scriptsize] at (\x,-0.1) {\lab};}
  \foreach \x/\n in {1.05/4, 1.4/8, 1.75/5, 2.1/3} {
    \foreach \k in {1,...,\n} { \fill[navy] (\x,{0.2*\k}) circle (0.07); }}
  \node[font=\scriptsize] at (2.8,-0.62) {Minutes waited};
  \node[font=\bfseries\small] at (2.8,2.15) {Line A};
  % ── Line B: widely spread, two customers stranded ──
  \begin{scope}[xshift=7.0cm]
    \draw[->, thick] (-0.2,0) -- (6.0,0);
    \foreach \x/\lab in {0/0, 0.7/2, 1.4/4, 2.1/6, 2.8/8, 3.5/10, 4.2/12, 4.9/14, 5.6/16} {
      \draw (\x,-0.07) -- (\x,0.07); \node[below, font=\scriptsize] at (\x,-0.1) {\lab};}
    \foreach \x/\n in {0.35/4, 0.7/3, 1.05/2, 1.4/2, 1.75/3, 2.1/2, 2.45/2, 3.85/1, 4.55/1} {
      \foreach \k in {1,...,\n} { \fill[goldacc] (\x,{0.2*\k}) circle (0.07); }}
    \node[font=\scriptsize] at (2.8,-0.62) {Minutes waited};
    \node[font=\bfseries\small] at (2.8,2.15) {Line B};
  \end{scope}
\end{tikzpicture}
\end{center}
\vspace{3pt}

\normalsize
\begin{enumerate}[label=\textbf{\alph*.}, leftmargin=1.8em, itemsep=4pt]
  \item How many customers were timed in Line A? \blank{1.6cm} \quad In Line B? \blank{1.6cm}

        \vspace{3pt}
        Counting in from both ends, the middle wait in Line A is \blank{1.6cm} minutes and the
        middle wait in Line B is \blank{1.6cm} minutes.
  \item Both lines have the same middle wait. \textbf{Which line would you rather join, and
        why?} Point at customers on the plots, not at the middle number.
        \answerspace{1.6cm}{}
  \item The manager reports one sentence to head office: \emph{``The typical wait is 4 minutes
        in both lines.''} Say exactly what that sentence hides about Line B.
        \answerspace{1.6cm}{}
  \item Two Line B customers sit far to the right of everybody else, at \textbf{11} and
        \textbf{13} minutes. What would make one of those a recording mistake? What would make
        it a real customer? Should it be left in the data either way?
        \answerspace{2.0cm}{}
\end{enumerate}
\end{scenariobox}

\boxguard[16]
\begin{scenariobox}[2. Three Pictures, Three Stories]{navy}
Three different things were measured; each is drawn as a histogram. \textbf{I} --- sale prices
of the 40 homes sold in one town last year. \textbf{II} --- scores of 32 students on an
end-of-unit quiz most of the class found easy. \textbf{III} --- heights of the 40 adults in a
mixed volleyball club.

\vspace{4pt}
\begin{center}
\begin{tikzpicture}[scale=0.78]
  % ── I: home prices, long right tail ──
  \foreach \i/\c in {0/6, 1/14, 2/10, 3/5, 4/3, 5/1, 6/1} {
    \fill[navy] ({\i*0.5},0) rectangle ({(\i+1)*0.5},{\c*0.14});
    \draw[white, very thin] ({\i*0.5},0) rectangle ({(\i+1)*0.5},{\c*0.14});}
  \draw[->, thick] (0,0) -- (0,2.5);
  \draw[->, thick] (0,0) -- (3.9,0);
  \foreach \y/\lab in {0/0, 0.7/5, 1.4/10} {
    \draw (-0.07,\y) -- (0.07,\y); \node[left, font=\scriptsize, xshift=-1pt] at (-0.07,\y) {\lab};}
  \foreach \x/\lab in {0/100, 1/300, 2/500, 3/700} {
    \draw (\x,-0.07) -- (\x,0.07); \node[below, font=\scriptsize] at (\x,-0.1) {\lab};}
  \node[font=\scriptsize] at (1.75,-0.68) {Sale price (\$1000s)};
  \node[font=\bfseries\small] at (1.75,2.75) {I};
  % ── II: quiz scores, long left tail ──
  \begin{scope}[xshift=4.9cm]
    \foreach \i/\c in {0/1, 1/2, 2/5, 3/11, 4/13} {
      \fill[goldacc] ({\i*0.7},0) rectangle ({(\i+1)*0.7},{\c*0.14});
      \draw[white, very thin] ({\i*0.7},0) rectangle ({(\i+1)*0.7},{\c*0.14});}
    \draw[->, thick] (0,0) -- (0,2.5);
    \draw[->, thick] (0,0) -- (3.9,0);
    \foreach \y/\lab in {0/0, 0.7/5, 1.4/10} {
      \draw (-0.07,\y) -- (0.07,\y); \node[left, font=\scriptsize, xshift=-1pt] at (-0.07,\y) {\lab};}
    \foreach \x/\lab in {0/50, 1.4/70, 2.8/90} {
      \draw (\x,-0.07) -- (\x,0.07); \node[below, font=\scriptsize] at (\x,-0.1) {\lab};}
    \node[font=\scriptsize] at (1.75,-0.68) {Quiz score};
    \node[font=\bfseries\small] at (1.75,2.75) {II};
  \end{scope}
  % ── III: club heights, symmetric with two humps and an empty stretch ──
  \begin{scope}[xshift=9.8cm]
    \foreach \i/\c in {0/2, 1/5, 2/9, 3/4, 4/0, 5/0, 6/4, 7/9, 8/5, 9/2} {
      \fill[navy] ({\i*0.35},0) rectangle ({(\i+1)*0.35},{\c*0.14});
      \draw[white, very thin] ({\i*0.35},0) rectangle ({(\i+1)*0.35},{\c*0.14});}
    \draw[->, thick] (0,0) -- (0,2.5);
    \draw[->, thick] (0,0) -- (3.9,0);
    \foreach \y/\lab in {0/0, 0.7/5, 1.4/10} {
      \draw (-0.07,\y) -- (0.07,\y); \node[left, font=\scriptsize, xshift=-1pt] at (-0.07,\y) {\lab};}
    \foreach \x/\lab in {0/58, 1.4/66, 2.8/74} {
      \draw (\x,-0.07) -- (\x,0.07); \node[below, font=\scriptsize] at (\x,-0.1) {\lab};}
    \node[font=\scriptsize] at (1.75,-0.68) {Height (inches)};
    \node[font=\bfseries\small] at (1.75,2.75) {III};
  \end{scope}
\end{tikzpicture}
\end{center}
\vspace{3pt}

\normalsize
\begin{enumerate}[label=\textbf{\alph*.}, leftmargin=1.8em, itemsep=4pt]
  \item For each picture: where do most of the values pile up, and which way do the stragglers
        run?

        \vspace{3pt}
        \renewcommand{\arraystretch}{1.2}
        \begin{tabularx}{\linewidth}{@{} c X X @{}}
        \rowcolor{sky}
        \textbf{Picture} & \textbf{Most of the values sit\ldots} &
        \textbf{The few stragglers run toward\ldots} \\
        \hline
        \textbf{I} & \blank{4.0cm} & \blank{4.0cm} \\
        \textbf{II} & \blank{4.0cm} & \blank{4.0cm} \\
        \textbf{III} & \blank{4.0cm} & \blank{4.0cm} \\
        \end{tabularx}
  \item Two of these three pictures lean, and they lean in \emph{opposite} directions. In your
        own words, describe how those two differ from each other.
        \answerspace{1.6cm}{}
  \item Picture III has two humps with an empty stretch between them. Suppose somebody told you
        only \emph{``The typical height in this club is about 67 inches.''} What would that tell
        you about the club --- and what would it leave out?
        \answerspace{1.8cm}{}
  \item A student summarized all three pictures by reporting only the average. Say what that
        summary misses in each case. Then write \textbf{one sentence} about Picture II that a
        person who had never seen it could still trust.
        \answerspace{2.2cm}{}
\end{enumerate}
\end{scenariobox}

% ── Part 2: QuickNotes (the debrief fills this) — target: half a page ────────
\boxguard[14]
\begin{tcolorbox}[enhanced, colback=sky, colframe=navy, boxrule=0.8pt, arc=2mm,
    left=3mm, right=3mm, top=1.5mm, bottom=1.5mm, breakable,
    fonttitle=\bfseries\small\color{white}, title={QuickNotes: Describing a Quantitative Distribution},
    attach boxed title to top left={yshift=-2mm, xshift=4mm},
    boxed title style={colback=navy, colframe=navy, arc=1.5mm}]
\small
\begin{itemize}[leftmargin=1.1em, itemsep=3pt]
  \item \textbf{Shape.} \textbf{Skewed right:} the \blank{1.6cm} tail is longer.
        \textbf{Skewed left:} the \blank{1.6cm} tail is longer. \textbf{Symmetric:} the left
        half is roughly a \blank{2.6cm} of the right.
        \textbf{The name follows the \blank{1.8cm}, not the \blank{1.6cm}.}
  \item \textbf{Peaks.} One main peak: \blank{2.2cm}. Two prominent peaks: \blank{2.2cm}.
        All bars about equal, no peak: \blank{2.4cm}.
  \item \textbf{Center} --- the \blank{2.2cm} value. \textbf{Variability (spread)} --- how far
        the values \blank{3.0cm}. Line A and Line B: same center, different spread.
  \item \textbf{Unusual features.} \textbf{Outlier:} a value unusually small or large
        \blank{3.4cm} the rest of the data. \textbf{Gap:} a stretch where \blank{2.6cm} was
        observed. \textbf{Cluster:} values concentrated together, set off by \blank{1.8cm}.
  \item \textbf{Four things or nothing:} \textbf{S}hape, \textbf{O}utliers, \textbf{C}enter,
        \textbf{S}pread --- remember it as \blank{1.8cm} --- and every one stated
        \blank{2.4cm}, naming the variable and its units.
  \item Describing \emph{these} data is not a claim about a larger \blank{2.4cm}; it can only
        suggest something worth \blank{2.4cm} later.
\end{itemize}
\end{tcolorbox}

% ── Part 3: the Application (worked together, ~6 min) ────────────────────────
\boxguard[14]
\begin{notesbox}{Application: The Sleep Survey}
We will do this one together. Thirty students reported their hours of sleep on a school night.

\vspace{3pt}
\begin{center}
\begin{tikzpicture}[scale=0.80]
  \draw[->, thick] (-0.2,0) -- (8.2,0);
  \foreach \x/\lab in {0.6/4, 1.8/5, 3.0/6, 4.2/7, 5.4/8, 6.6/9, 7.8/10} {
    \draw (\x,-0.07) -- (\x,0.07); \node[below, font=\scriptsize] at (\x,-0.1) {\lab};}
  \foreach \x/\n in {0.6/1, 3.0/1, 3.6/2, 4.2/4, 4.8/5, 5.4/7, 6.0/6, 6.6/3, 7.2/1} {
    \foreach \k in {1,...,\n} { \fill[navy] (\x,{0.2*\k}) circle (0.07); }}
  \node[font=\scriptsize] at (4.0,-0.62) {Hours of sleep on a school night};
\end{tikzpicture}
\end{center}
\vspace{2pt}

\begin{enumerate}[label=\textbf{\alph*.}, leftmargin=1.8em, itemsep=4pt]
  \item How many students are shown? \blank{1.8cm} \quad One value sits apart from all the
        others: \blank{2.0cm} hours.
  \item Which way does this picture lean? \blank{3.0cm} \quad Say what that longer tail means
        \emph{for these students} --- not for the graph.
        \answerspace{1.6cm}{}
  \item Now write the complete description in one paragraph: all four parts, in context, with
        the variable and its units named.
        \answerspace{2.4cm}{}
\end{enumerate}
\end{notesbox}

\end{document}
